%% parts/p14/master-tables.tex --- Part 14 apparatus
\chapter{Master Cross-Reference Tables, Glossary, and Indexes}
\label{sec:p14-master-tables}

These tables are retrieval apparatus. A row abbreviates a theorem or definition
at the cited destination and never strengthens it. In particular, complexity
rows retain the decision problem and input convention, strength rows retain
their base theories, and correspondence rows retain both the frame and
algebraic hypotheses.

\section{Logic and decision-problem table}
\label{sec:p14-decision-table}

\begin{longtable}{@{}p{0.20\textwidth}p{0.23\textwidth}p{0.20\textwidth}p{0.27\textwidth}@{}}
\toprule
Language or system & Problem & Status & Qualification \\
\midrule
\endfirsthead
\toprule
Language or system & Problem & Status & Qualification \\
\midrule
\endhead
CPL & satisfiability & NP-complete & polynomial encodings;
Theorem~\ref{thm:p04-cook-levin} \\
Horn CPL & satisfiability & P-complete & forward chaining \cite{DowlingGallier1984} \\
2-CNF & satisfiability & NL-complete & implication graph \cite{AspvallPlassTarjan1979} \\
QBF & truth & PSPACE-complete & fully quantified formulas \cite{Stockmeyer1976} \\
first-order logic & validity & r.e.-complete, undecidable & arbitrary structures \cite{Church1936,Turing1937} \\
first-order logic & finite validity & not r.e. & binary-relation encodings \cite{Trakhtenbrot1950} \\
monadic first-order logic & satisfiability & NEXPTIME-complete & closed formulas
over unary relation symbols; no identity, constants, or function symbols;
Table~\ref{tab:p02-decidable-fragments} \cite{Lewis1980} \\
two-variable first-order logic & satisfiability & NEXPTIME-complete & equality allowed \cite{GradelKolaitisVardi1997} \\
guarded fragment & satisfiability & decidable & atom-guarded quantification \cite{Gradel1999} \\
\(\mathcal{ALC}\) with general TBoxes & satisfiability & EXPTIME-complete & Definition~\ref{def:p09-alc}; \cite{BaaderHorrocksLutzSattler2017} \\
modal \(\mathsf K\) & satisfiability & PSPACE-complete & arbitrary frames \cite{Ladner1977} \\
modal \(\mathsf{S4}\) & satisfiability & PSPACE-complete & reflexive transitive frames \cite{Ladner1977} \\
modal \(\mathsf{S5}\) & satisfiability & NP-complete & equivalence frames \cite{Ladner1977} \\
LTL & satisfiability & PSPACE-complete & infinite words \cite{SistlaClarke1985} \\
CTL & model checking & P-complete & finite explicit total systems;
Theorem~\ref{thm:p09-ctl-star-complexity} \\
CTL & satisfiability & EXPTIME-complete & exact CTL state-formula language;
Theorem~\ref{thm:p09-ctl-star-complexity} \\
CTL* & model checking & PSPACE-complete & formula and finite explicit total
system are both input; Theorem~\ref{thm:p09-ctl-star-complexity} \\
CTL* & satisfiability & 2EXPTIME-complete & state/path typing retained;
Theorem~\ref{thm:p09-ctl-star-complexity} \\
PDL & satisfiability & EXPTIME-complete & regular programs;
Theorem~\ref{thm:p09-pdl-decision} \\
PCTL & model checking & polynomial time & finite rational DTMC or MDP,
binary rational data, unary bounded-until integers, explicit min/max MDP
operator; Theorem~\ref{thm:p09-pctl-model-checking} \\
CSL & model checking & effective with certified numerical boundary & finite
rational CTMC; time-bounded thresholds require rigorous error intervals and
equality requires separate exact analysis; Theorem~\ref{thm:p09-csl-model-checking} \\
propositional likelihood logic & satisfiability & NP-complete & Boolean
combinations of binary-encoded rational linear inequalities over event
probabilities; Theorem~\ref{thm:p09-likelihood-logic} \\
modal \(\mu\)-calculus & satisfiability & EXPTIME-complete & exact positive
fixed-point language; \hyperref[entry:p08-modal-mu]{Part VIII system card}
\cite{Kozen1983} \\
modal \(\mu\)-calculus & finite-state model checking & polynomial-time
equivalent to parity games; P-hard; in
\(\mathsf{UP}\cap\mathsf{coUP}\); quasipolynomial-time; membership in P open &
finite explicitly represented transition systems;
\hyperref[entry:p08-modal-mu]{Part VIII system card} and
Table~\ref{tab:p08-temporal-comparison}
\cite{EmersonJutla1999,CaludeJainKhoussainovLiStephan2017} \\
Presburger arithmetic & truth & decidable & addition/order; lower bound in \cite{FischerRabin1974} \\
true first-order arithmetic & truth & undecidable, non-r.e. & standard naturals \cite{Shoenfield1967} \\
finite-template CSP & decision & P or NP-complete & fixed template \cite{Zhuk2020} \\
finite normal programs & stable-model existence & NP-complete & propositional,
nondisjunctive; Theorem~\ref{thm:p09-stable-complexity} \\
finite normal programs & credulous atom consequence & NP-complete &
propositional, nondisjunctive; Theorem~\ref{thm:p09-stable-complexity} \\
finite normal programs & skeptical atom consequence & coNP-complete &
propositional, nondisjunctive; Theorem~\ref{thm:p09-stable-complexity} \\
conjunctive queries & fixed-query evaluation & \(\mathsf{AC}^{0}\) & finite
equality-free relational data; Theorem~\ref{thm:p09-cq-data-evaluation} \\
conjunctive queries & combined evaluation & NP-complete & query and finite
database are input; Theorem~\ref{thm:p09-cq-combined-evaluation} \\
conjunctive queries & containment & NP-complete & two same-arity equality-free
queries; Theorem~\ref{thm:p09-cq-containment} \\
fixed Datalog program and query & query evaluation & in P, with some fixed
query P-complete & finite extensional database is the input;
Theorem~\ref{thm:p09-datalog-complexity} \\
Datalog & combined query evaluation & EXPTIME-complete & program, query, and
finite database are all input; Theorem~\ref{thm:p09-datalog-complexity} \\
\bottomrule
\end{longtable}

The table distinguishes taxonomic, data, query, combined, and formula
complexity. ``Decidable''
without a class is intentionally weaker than a complete complexity
classification, and ``undecidable'' does not identify an arithmetical degree
unless one is stated at the source theorem.

\section{Consistency and proof-theoretic strength}
\label{sec:p14-strength-table}

There is no single unconditional total ordering containing every foundational
theory. The following chains use the indicated interpretation or proof-theoretic
comparison; arrows are not set inclusion unless the row says so.

Table~\ref{tab:p06-consistency-strength-master} and
Figure~\ref{fig:p06-consistency-strength} are the master strength surfaces.
They keep the arithmetic and reverse-mathematics tracks separate, state the
PA/ACA-zero conservation bridge, distinguish set-theory extension from
equiconsistency, and label every large-cardinal edge conditionally.  An ordinal
analysis, an interpretation, a conservation theorem, and a relative
consistency theorem are four different comparison relations; neither master
surface converts one into another without an explicit theorem.

\section{Expressive-power map}
\label{sec:p14-expressive-table}

\begin{longtable}{@{}p{0.21\textwidth}p{0.27\textwidth}p{0.42\textwidth}@{}}
\toprule
Comparison & Inclusion or equivalence & Witnessing boundary \\
\midrule
\endfirsthead
\toprule
Comparison & Inclusion or equivalence & Witnessing boundary \\
\midrule
\endhead
modal logic vs first-order logic & standard translation embeds modal logic in FO & bisimulation-invariant boundary \cite{VanBenthem1983} \\
CPL vs first-order logic & propositional formulas embed as nullary-predicate cases & domain-sensitive quantification \cite{Enderton2001} \\
FO vs existential SO on finite structures & proper inclusion & ESO captures NP \cite{Fagin1974} \\
FO vs least fixed-point logic on ordered finite structures & proper inclusion & LFP captures P \cite{Immerman1986} \\
LFP vs ESO on ordered finite structures & inclusion if \(\mathsf P\subseteq\mathsf{NP}\) & captured-class equality boundary \cite{Immerman1999} \\
LTL vs FO on words & LTL with past captures FO point properties & Kamp theorem \cite{Kamp1968} \\
CTL vs LTL & incomparable & branching/path witnesses \cite{EmersonHalpern1986} \\
CTL and LTL vs CTL* & both proper fragments & unrestricted state/path nesting \cite{EmersonHalpern1986} \\
modal \(\mu\)-calculus vs modal logic & proper extension & fixed-point reachability \cite{Kozen1983} \\
full SO vs Henkin SO & same syntax, different validities & full versus general relation ranges \cite{Henkin1950,Shapiro1991} \\
\bottomrule
\end{longtable}

All comparisons hold over the displayed structure class and semantics. A
translation that preserves satisfiability need not preserve formula size,
validity, proof length, or intended interpretation.

\section{Modal axiom--frame--algebra correspondences}
\label{sec:p14-modal-correspondence}

The rows use the normal-modal correspondence and complex-algebra translations
in \cite{BlackburnDeRijkeVenema2001,Goldblatt1989}.

\begin{tabularx}{\textwidth}{@{}l l l X@{}}
\toprule
Axiom & Frame condition & Algebraic inequality & System use \\
\midrule
\(\Box(p\to q)\to(\Box p\to\Box q)\) & none beyond normality & \(\Box\) preserves finite meets and \(1\) & \(\mathsf K\) \\
\(\Box p\to p\) & reflexive & \(\Box a\le a\) & \(\mathsf T\) \\
\(\Box p\to\Box\Box p\) & transitive & \(\Box a\le\Box\Box a\) & axiom 4 \\
\(p\to\Box\Diamond p\) & symmetric & \(a\le\Box\Diamond a\) & axiom B \\
\(\Box p\to\Diamond p\) & serial & \(\Box0=0\) & axiom D \\
\(\Diamond p\to\Box\Diamond p\) & Euclidean & \(\Diamond a\le\Box\Diamond a\) & axiom 5 \\
T + 4 & reflexive, transitive & interior algebra & \(\mathsf{S4}\) \\
T + 5 & equivalence relation & monadic Boolean algebra & \(\mathsf{S5}\) \\
\bottomrule
\end{tabularx}

The table records standard correspondence for normal propositional modal
logic. It does not transfer unchanged to neighborhood, many-valued, first-order
modal, dynamic, or nonnormal semantics.

\section{Glossary}
\label{sec:p14-glossary}

\begin{description}[style=nextline,leftmargin=0pt]
\item[admissible rule] A rule whose conclusion is derivable after every
substitution making all premises derivable, whether or not the rule itself is
derivable.
\item[algebraic semantics] An interpretation of logical operations and
consequence in a declared class of algebras; see Chapter~\ref{sec:p10-algebraic-and-categorical-logic}.
\item[compactness] The property that semantic consequence from a set is already
consequence from a finite subset; first-order compactness is
Theorem~\ref{thm:fol-compactness}.
\item[complete theory] A consistent theory deciding every sentence of its
language; this differs from completeness of a calculus.
\item[congruence] An operation-compatible equivalence relation on an algebra;
see Definition~\ref{def:p10-algebra}.
\item[conservative extension] An extension proving no new sentences in the old
language beyond those already provable in the base theory.
\item[decidable set] A set with a total computable characteristic function.
\item[derivability] Existence of a derivation in a named calculus; never an
unsubscripted synonym for semantic consequence.
\item[effective procedure] A finitely specified method analyzed through a
declared formal model of computation.
\item[elementary embedding] An embedding preserving every first-order formula
with parameters.
\item[equisatisfiable] Having a model simultaneously in the sense that one
formula is satisfiable iff the other is; weaker than equivalence.
\item[filter] An upward-closed, finite-meet-closed subset of a lattice; see
Definition~\ref{def:p10-boolean-algebra}.
\item[finite validity] Truth in every finite structure of the fixed signature;
see Definition~\ref{def:p09-finite-validity}.
\item[frame] The relational component of a Kripke model, before a valuation is
added.
\item[free algebra] An algebra satisfying the extension universal property from
a chosen generating set.
\item[homomorphism] A map preserving every operation or relation specified by
the context.
\item[interpretation of theories] A translation of one theory into another
with a declared domain formula and preservation of axioms/proofs.
\item[invariant] A property unchanged under a specified class of transformations;
the transformation class is part of the assertion.
\item[Kripke model] A frame together with a valuation or domain/interpretation
data appropriate to the modal language.
\item[Lindenbaum--Tarski algebra] The formula algebra modulo provable
interderivability; see Definition~\ref{def:p10-lindenbaum-algebra}.
\item[model checking] Determining whether a fixed finite structure/state
satisfies a formula; distinct from satisfiability over all structures.
\item[normal form] A formula satisfying a declared syntactic restriction and
connected to the input by equivalence or equisatisfiability as explicitly stated.
\item[proof sketch] An incomplete proof presentation whose omitted steps remain
owed and which is never cited as a complete proof.
\item[realizability] An interpretation assigning computational witnesses to
formulas under a specified applicative structure.
\item[relative consistency] A metatheoretic implication
\(\operatorname{Con}(T)\to\operatorname{Con}(U)\), not an absolute consistency proof.
\item[saturation] Realization of specified types over parameter sets below a
declared cardinal bound.
\item[semantic consequence] Truth preservation over a declared class of
interpretations; see Definition~\ref{def:p11-model-consequence}.
\item[soundness] The assertion that derivability in a named calculus implies
semantic consequence under a named semantics.
\item[stable model] A fixed point of the Gelfond--Lifschitz reduct construction;
see Definition~\ref{def:p09-stable-model}.
\item[Stone space] The compact zero-dimensional Hausdorff dual of a Boolean
algebra; see Definition~\ref{def:p10-stone-space}.
\item[substitution] Capture-avoiding replacement with the exact term/formula
and free-variable conditions stated.
\item[theory] A set of sentences, often deductively closed only when explicitly
said so.
\item[type] A consistent set of formulas in specified free variables over a
declared parameter set.
\item[ultrafilter] A maximal proper filter; see Definition~\ref{def:p10-boolean-algebra}.
\item[validity] Truth in every interpretation in a declared class.
\item[variety] A class of similar algebras axiomatized by equations;
equivalently an HSP-closed class when arbitrary products include the empty
product, by Theorem~\ref{thm:p10-hsp}.
\end{description}

\section{Named-theorem index}
\label{sec:p14-theorem-index}

\begin{longtable}{@{}p{0.35\textwidth}p{0.10\textwidth}p{0.13\textwidth}p{0.30\textwidth}@{}}
\toprule
Result & Kind & Statement & Proof/detail \\
\midrule
\endfirsthead
\toprule
Result & Kind & Statement & Proof/detail \\
\midrule
\endhead
\hbinputtablerows{generated/named-theorem-index.tex}
\bottomrule
\end{longtable}

\section{Symbol, person, and subject indexes}
\label{sec:p14-index-surfaces}

\subsection*{Symbol index}
\begin{longtable}{@{}p{0.12\textwidth}p{0.19\textwidth}p{0.10\textwidth}p{0.11\textwidth}p{0.09\textwidth}p{0.27\textwidth}@{}}
\toprule
Symbol & Role & Scope & Registration & First use & Principal uses / notes \\
\midrule
\endfirsthead
\toprule
Symbol & Role & Scope & Registration & First use & Principal uses / notes \\
\midrule
\endhead
\hbinputtablerows{generated/symbol-index.tex}
\bottomrule
\end{longtable}

The generated registry in Chapter~\ref{sec:p01-symbol-registry} is the symbol
index: every row records the printed symbol, spelling, role, scope, first-use
locator, registration-page destination, and principal-use notes; the compact
alphabetical projection above links back to each registration row. The capsule sequence in
Section~\ref{sec:p13-biographies} is historical, not alphabetical; surname-
sorted person entries in the final unified index point to that section. The
named-theorem table above is generated from every live named theorem and lemma
destination, with the statement page and the recorded proof/annex/citation
detail.
The backmatter index is the unified subject index generated from the exact
inherited entries together with editorial entries. Thus the four retrieval
surfaces are distinct without printing four duplicate alphabetical books.

\section{Further-reading map}
\label{sec:p14-further-reading}

\begin{tabularx}{\textwidth}{@{}l X@{}}
\toprule
Part & Starting sources \\
\midrule
I--II & \cite{Enderton2001,ChiswellHodges2007,OpenLogic2026} \\
III & \cite{Hodges1993,Marker2002,Marker2024} \\
IV & \cite{BoolosBurgessJeffrey2007,Shoenfield1967,AroraBarak2009} \\
V & \cite{Kaye1991,Boolos1993,Simpson2009} \\
VI & \cite{Jech2003,Kunen2011,Kanamori2003} \\
VII & \cite{TroelstraSchwichtenberg2000,Prawitz1965,GirardLafontTaylor1989} \\
VIII & \cite{BlackburnDeRijkeVenema2001,ChagrovZakharyaschev1997,Hajek1998} \\
IX & \cite{Libkin2004,BaaderHorrocksLutzSattler2017,ClarkeGrumbergPeled1999} \\
X & \cite{BurrisSankappanavar1981,Johnstone1982,MacLaneMoerdijk1992} \\
XI & \cite{Halbach2011,Keefe2000,BeallRestall2006} \\
XII & \cite{RobinsonVoronkov2001,SMTLIB2025,LeanManual2026} \\
XIII--XIV & primary works in Section~\ref{sec:p13-chronology} and the unified bibliography \\
\bottomrule
\end{tabularx}
